\documentclass[12pt,letterpaper]{article}
\usepackage[utf8]{inputenc}
\usepackage[spanish,mexico]{babel}
\usepackage[T1]{fontenc}
\usepackage{geometry}
\usepackage{parskip}
\usepackage{hyperref}

\geometry{margin=2.5cm}
\hypersetup{colorlinks=true,linkcolor=blue,citecolor=blue,urlcolor=blue}

\title{\'Etica profesional y pedagog\'{\i}a del bien com\'{u}n:\\ fundamentos para la pr\'actica tutorial en entornos presenciales y virtuales}
\author{---}
\date{}

\begin{document}

\maketitle

\begin{abstract}
La \'etica profesional del docente no puede ejercerse sin una comprensi\'on previa del concepto mismo de \'etica. Este documento explora las ra\'{\i}ces del concepto \'etico a partir de la obra de Roger-Pol Droit, lo vincula con los marcos de la \'etica profesional docente y lo proyecta hacia la pedagog\'{\i}a del bien com\'un propuesta por la UNESCO (2015). Se presentan los hallazgos principales de esta corriente, su relaci\'on con la \'etica, y se analiza su aplicabilidad en instituciones educativas --- incluyendo entornos remotos y virtuales --- donde la mediaci\'on tecnol\'ogica replantea las preguntas \'eticas fundamentales.
\end{abstract}

\section{Introducci\'on}

La palabra \textit{\'etica} se ha vuelto omnipresente en el discurso educativo contempor\'aneo, pero su uso frecuente no implica una comprensi\'on profunda. Como se\~nala Roger-Pol Droit (2010), pocas veces se reflexiona sobre su verdadero significado m\'as all\'a del sentido com\'un. Para un profesional de la educaci\'on que aspira a acompa\~nar \'eticamente a sus estudiantes, comprender qu\'e es la \'etica --- de d\'onde viene, c\'omo opera, qu\'e exige --- no es un lujo te\'orico sino una necesidad pr\'actica. Este documento parte de esa premisa: no se puede hablar de \'etica profesional sin haber entendido primero qu\'e es la \'etica.

A partir de esa comprensi\'on fundamental, se exploran dos dimensiones complementarias. La primera es la \'etica profesional del docente, que traduce el concepto abstracto de \'etica en principios concretos de acci\'on en el aula y en la tutor\'{\i}a. La segunda es la pedagog\'{\i}a del bien com\'un, un marco propuesto por la UNESCO (2015) que sit\'ua la educaci\'on como un compromiso colectivo y no solo como un servicio individual. Ambas dimensiones convergen en una pregunta pr\'actica: \textquestiondown c\'omo forma el tutor a sus estudiantes para una ciudadan\'{\i}a \'etica cuando las condiciones de aprendizaje son cada vez m\'as diversas, incluyendo entornos remotos y virtuales?

\section{El concepto de \'etica: de Droit al aula}

\subsection{Las aventuras de una palabra}

Droit (2010) rastrea el origen de la palabra \textit{\'etica} hasta el vocablo griego \textit{\`ethos}, cuyo significado era m\'as amplio que el actual. \textit{Ethos} pod\'{\i}a referirse al h\'abitat de una especie, al car\'acter de una persona o a las costumbres de una sociedad. De all\'{\i} surgi\'o \textit{\`ethik\'e}, t\'ermino utilizado por Arist\'oteles para designar el conocimiento relativo a los comportamientos humanos. La \'etica naci\'o, entonces, como una reflexi\'on sobre la conducta, no como un c\'odigo fijo.

Una distinci\'on clave que Droit establece es la diferencia entre \textit{\'etica} y \textit{moral}. Originalmente sin\'onimos (los romanos tradujeron \textit{\`ethik\'e} como \textit{moralia}), ambos t\'erminos se separaron en la modernidad. La moral pas\'o a asociarse con normas y valores heredados de la tradici\'on, la cultura o la religi\'on; la \'etica, en cambio, se relacion\'o con el trabajo reflexivo necesario para responder a situaciones nuevas para las cuales no existen reglas predefinidas. Desde esta perspectiva, la \'etica puede entenderse como una moral en construcci\'on --- una deliberaci\'on viva, no un manual muerto.

\subsection{Implicaciones para la formaci\'on profesional}

Esta distinci\'on tiene consecuencias directas para la pr\'actica docente. Como se\~nalan Ramos Serpa y L\'opez Falc\'on (2019), la formaci\'on \'etica del profesional no puede limitarse a la transmisi\'on de c\'odigos deontol\'ogicos; debe desarrollar la capacidad de analizar situaciones complejas, reconocer conflictos de valores y fundamentar racionalmente las decisiones. El docente \'etico no es quien memoriza un reglamento, sino quien sabe deliberar cuando el reglamento no cubre el caso concreto.

La \'etica profesional, entonces, no es un adorno ni una burocracia. Es la competencia de reconocer un dilema donde otros solo ven una instrucci\'on, y de decidir con criterio donde no hay una respuesta \'unica. Esta capacidad de deliberaci\'on pr\'actica --- lo que Arist\'oteles llam\'o \textit{phr\'onesis} --- constituye el n\'ucleo de la \'etica profesional docente y el puente hacia la pedagog\'{\i}a del bien com\'un.

\section{La pedagog\'{\i}a del bien com\'un}

\subsection{Origen y definici\'on}

En 2015, la UNESCO public\'o el informe \textit{Replantear la educaci\'on: \textquestiondown Hacia un bien com\'un mundial?}, coordinado por Sobhi Tawil y Rita Locatelli. El documento propone un cambio de paradigma: entender la educaci\'on no solo como un bien p\'ublico (responsabilidad del Estado) ni como un bien privado (servicio individual), sino como un \textbf{bien com\'un}. Esto implica que la educaci\'on es un compromiso colectivo que nos sostiene a todos y cuya gobernanza debe ser participativa, inclusiva y orientada al bienestar de la comunidad (UNESCO, 2015).

Locatelli (2018) profundiza esta idea al distinguir entre lo p\'ublico (gestionado por el Estado) y lo com\'un (co-gestionado por la comunidad). En una sociedad plural y cambiante, el Estado no puede ni debe ser el \'unico garante de la educaci\'on; la sociedad civil, las comunidades locales y los propios educandos deben participar activamente en la definici\'on de sus fines y medios.

\subsection{Principios fundamentales}

A partir del an\'alisis de la literatura acad\'emica reciente (UNESCO, 2015; Locatelli, 2018; Merino Fern\'andez, 2022; Croda et al., 2022; Urrego Estrada et al., 2022), pueden identificarse varios principios que definen la pedagog\'{\i}a del bien com\'un:

\begin{enumerate}
  \item \textbf{Corresponsabilidad.} La educaci\'on es tarea de todos, no solo del Estado ni de la instituci\'on escolar. La comunidad participa en su definici\'on y sostenimiento (Locatelli, 2018; Urrego Estrada et al., 2022).
  \item \textbf{Inclusi\'on y equidad.} El bien com\'un no puede lograrse si algunos sectores quedan excluidos. La pedagog\'{\i}a del bien com\'un exige eliminar las barreras que impiden la participaci\'on plena (UNESCO, 2015; Flores Crespo, 2018).
  \item \textbf{Sostenibilidad.} La educaci\'on como bien com\'un debe pensarse a largo plazo, en t\'erminos intergeneracionales. No solo importa el aprendizaje de hoy, sino la capacidad de las generaciones futuras para aprender y participar (UNESCO, 2015).
  \item \textbf{Di\'alogo y participaci\'on democr\'atica.} La definici\'on de lo que es valioso ense\~nar no puede ser impuesta unilateralmente. La pedagog\'{\i}a del bien com\'un implica la construcci\'on colectiva de acuerdos sobre los fines de la educaci\'on (Merino Fern\'andez, 2022; Croda et al., 2022).
  \item \textbf{Trascendencia del individuo.} El aprendizaje no es solo un logro personal; es un bien que se multiplica cuando se comparte. Un estudiante educado no solo mejora su propia vida, sino que contribuye al bienestar de su comunidad (Gracia-Caland\'{\i}n y Tamarit-L\'opez, 2021).
\end{enumerate}

\section{Relaci\'on entre la \'etica y la pedagog\'{\i}a del bien com\'un}

La conexi\'on entre ambos conceptos no es accidental; es estructural. La \'etica, entendida como reflexi\'on sobre la conducta y deliberaci\'on ante lo nuevo, encuentra en la pedagog\'{\i}a del bien com\'un un campo de aplicaci\'on concreto. Inversamente, la pedagog\'{\i}a del bien com\'un necesita de la \'etica como fundamento para no convertirse en una consigna vac\'{\i}a.

Tres puntos de convergencia merecen destacarse:

\begin{enumerate}
  \item \textbf{El bien com\'un como criterio \'etico.} La \'etica profesional del tutor no puede limitarse al bien individual del estudiante. Como se\~nala Cortina (1986), las profesiones tienen bienes internos que las orientan; para la docencia, ese bien interno no es solo el aprendizaje individual, sino la formaci\'on de ciudadanos que sostengan el bien com\'un. Una decisi\'on \'etica que beneficia a un estudiante pero da\~na a la comunidad no es \'eticamente completa.
  \item \textbf{La deliberaci\'on como m\'etodo compartido.} Tanto la \'etica (Droit, 2010) como la pedagog\'{\i}a del bien com\'un (UNESCO, 2015) exigen di\'alogo y construcci\'on colectiva. No hay recetas; hay principios que se aplican mediante la deliberaci\'on informada. El docente que forma para el bien com\'un est\'a formando, al mismo tiempo, la capacidad \'etica de sus estudiantes.
  \item \textbf{La dignidad como l\'{\i}mite y horizonte.} Tanto Kant (desde la \'etica del deber) como Freire (desde la pedagog\'{\i}a de la liberaci\'on) coinciden en un punto: la persona es un fin, no un medio. La pedagog\'{\i}a del bien com\'un a\~nade que ese fin no es solo individual sino colectivo: cada persona realiza su dignidad en comunidad con otras.
\end{enumerate}

\section{Vinculaci\'on con la realidad institucional educativa}

\subsection{En entornos presenciales}

En la instituci\'on educativa presencial, la pedagog\'{\i}a del bien com\'un se traduce en pr\'acticas concretas: proyectos colaborativos que vinculan la escuela con la comunidad, asambleas estudiantiles donde se discuten las normas, tutor\'{\i}as que no solo atienden el rendimiento sino la integraci\'on social, y evaluaciones que consideran el contexto del estudiante sin renunciar al rigor acad\'emico.

El tutor, en este marco, no es un gestor de casos individuales sino un facilitador del bien com\'un. Su tarea \'etica incluye velar por que ning\'un estudiante quede excluido, que los recursos se distribuyan seg\'un necesidad (equidad, no igualdad ciega) y que las decisiones se tomen con transparencia y participaci\'on.

\subsection{En entornos remotos y virtuales}

La mediaci\'on tecnol\'ogica plantea desaf\'{\i}os espec\'{\i}ficos que la pedagog\'{\i}a del bien com\'un ayuda a abordar:

\begin{enumerate}
  \item \textbf{Brecha digital como exclusi\'on del bien com\'un.} Cuando un estudiante no tiene acceso a internet o a un dispositivo adecuado, no est\'a solo en desventaja individual: est\'a excluido del bien com\'un educativo. La responsabilidad \'etica del tutor incluye identificar estas brechas y buscar v\'{\i}as institucionales para cerrarlas, no solo adaptar la evaluaci\'on.

  \item \textbf{Privacidad y vigilancia.} En entornos virtuales, la frontera entre acompa\~nar y vigilar se vuelve delgada. Rivera Piragauta y Minelli de Oliveira (2017, 2020) advierten sobre el problema \'etico de la identidad digital y el plagio, pero tambi\'en sobre la tendencia a monitorear excesivamente a los estudiantes. La pedagog\'{\i}a del bien com\'un ofrece un criterio: las decisiones sobre vigilancia tecnol\'ogica deben tomarse con participaci\'on de la comunidad educativa y en funci\'on del bien colectivo, no de la eficiencia administrativa.

  \item \textbf{Aislamiento y comunidad virtual.} El aprendizaje remoto puede fragmentar el tejido comunitario que la educaci\'on presencial sostiene de forma natural. La \'etica del tutor en entornos virtuales incluye la tarea expl\'{\i}cita de construir comunidad, de asegurar que los estudiantes no est\'en solos frente a la pantalla. Esto implica dise\~nar espacios de interacci\'on genuina, no solo de transmisi\'on de contenidos.

  \item \textbf{Sesgo algor\'{\i}tmico y justicia.} Herramientas de inteligencia artificial que revisan tareas, predicen deserci\'on o clasifican estudiantes pueden incorporar sesgos que perpet\'uan desigualdades (O'Neil, 2017). La pedagog\'{\i}a del bien com\'un exige transparencia algor\'{\i}tmica y participaci\'on de la comunidad en la decisi\'on de qu\'e herramientas usar y bajo qu\'e criterios.
\end{enumerate}

\subsection{Hacia una \'etica tutorial para entornos mixtos}

La integraci\'on de lo presencial y lo virtual no es una opci\'on t\'ecnica sino \'etica. El tutor que trabaja en entornos mixtos necesita principios que operen en ambos contextos: la confidencialidad con l\'{\i}mites, la justicia como equidad, el respeto a la autonom\'{\i}a creciente del estudiante y la responsabilidad profesional. Estos principios, derivados de la \'etica profesional dialogada con la pedagog\'{\i}a del bien com\'un, ofrecen un marco portable que no depende de la plataforma ni de la modalidad.

\section{Conclusiones}

La \'etica profesional del docente no es un protocolo que se aplica, sino una capacidad que se cultiva. Comprender el concepto de \'etica --- su origen, su distinci\'on de la moral, su exigencia de deliberaci\'on --- es el primer paso para ejercerla. La pedagog\'{\i}a del bien com\'un a\~nade una dimensi\'on colectiva que impide reducir la \'etica docente al cuidado individual: el tutor no solo responde ante cada estudiante, sino ante la comunidad que habita.

En entornos remotos y virtuales, estas preguntas se vuelven m\'as urgentes. La tecnolog\'{\i}a no es \'eticamente neutra; reproduce o transforma las condiciones en que el bien com\'un se construye. El tutor \'etico del siglo XXI es quien sabe deliberar con criterio, cuidar sin invadir, y educar para que el bien aprendido se multiplique en comunidad.

\begin{thebibliography}{99}

\bibitem[UNESCO(2015)]{UNESCO2015}
UNESCO. (2015). \emph{Replantear la educaci\'on: \textquestiondown Hacia un bien com\'un mundial?} Par\'{\i}s: UNESCO. \url{https://unesdoc.unesco.org/ark:/48223/pf0000232697}

\bibitem[Locatelli(2018)]{Locatelli2018}
Locatelli, R. (2018). Education as a public and common good: Reframing the governance of education in a changing context. \emph{UNESCO Education Research and Foresight: Working Papers}, 22. \url{https://unesdoc.unesco.org/ark:/48223/pf0000261614}

\bibitem[Droit(2010)]{Droit2010}
Droit, R.-P. (2010). \emph{La \'etica explicada a todo el mundo}. Barcelona: Paid\'os.

\bibitem[Ramos Serpa y L\'opez Falc\'on(2019)]{Ramos2019}
Ramos Serpa, G. y L\'opez Falc\'on, A. (2019). Formaci\'on \'etica del profesional y \'etica profesional del docente. \emph{Estudios Pedag\'ogicos (Valdivia)}, 45(3), 185--199. \url{https://www.scielo.cl/scielo.php?pid=S0718-07052019000300185\&script=sci_arttext}

\bibitem[Flores Crespo(2018)]{Flores2018}
Flores Crespo, P. (2018). \textquestiondown La educaci\'on como bien com\'un? \emph{Metodolog\'{\i}a y Evaluaci\'on de la Calidad Educativa}, 20, 83--98. \url{https://doi.org/10.24310/metyper.2018.v0i20.4826}

\bibitem[Gracia-Caland\'{\i}n y Tamarit-L\'opez(2021)]{Gracia2021}
Gracia-Caland\'{\i}n, J. y Tamarit-L\'opez, I. (2021). Education as a common good from the capability approach. \emph{Journal of Philosophy of Education}, 55(5-6), 817--828. \url{https://doi.org/10.1111/1467-9752.12575}

\bibitem[Merino Fern\'andez(2022)]{Merino2022}
Merino Fern\'andez, J. V. (2022). La pedagog\'{\i}a del bien com\'un: \textquestiondown enfoque o teor\'{\i}a pedag\'ogica? \emph{A\&H, Revista de Artes, Humanidades y Ciencias Sociales}, N\'umero Especial, 26--118. \url{https://revistas.upaep.mx/index.php/ayh/article/view/336}

\bibitem[Croda, Ba\~nuelos y Lechuga(2022)]{Croda2022}
Croda, G., Ba\~nuelos, F. y Lechuga, G. (2022). Principios de la pedagog\'{\i}a del bien com\'un: experiencia de actores clave de la comunidad educativa. \emph{A\&H, Revista de Artes, Humanidades y Ciencias Sociales}, N\'umero Especial, 177--204. \url{https://revistas.upaep.mx/index.php/ayh/article/view/341}

\bibitem[Urrego Estrada et al.(2022)]{Urrego2022}
Urrego Estrada, G. A., Colorado Rend\'on, S. E. y Betancur Hern\'andez, L. F. (2022). Los bienes comunes desde una pedagog\'{\i}a social: un horizonte posible para pensar lo comunitario. \emph{Collectivus. Revista de Ciencias Sociales}, 9(2), 463--516. \url{https://doi.org/10.15648/Collectivus.vol9num2.2022.3530}

\bibitem[Rivera Piragauta y Minelli de Oliveira(2017)]{Rivera2017}
Rivera Piragauta, J. A. y Minelli de Oliveira, J. (2017). El problema \'etico de la identidad digital en la educaci\'on virtual. \emph{Revista Ibero-Americana de Educaci\'on}, 75(2), 41--58. \url{https://doi.org/10.35362/rie7522633}

\bibitem[Rivera Piragauta y Minelli de Oliveira(2020)]{Rivera2020}
Rivera Piragauta, J. A. y Minelli de Oliveira, J. (2020). M\'as all\'a del plagio: relevancia de la \'etica en entornos virtuales de aprendizaje. \emph{ETIC@NET}, 20(1), 156--185. \url{https://doi.org/10.30827/eticanet.v20i1.15140}

\bibitem[Parra Castrill\'on(2021)]{Parra2021}
Parra Castrill\'on, E. (2021). An\'alisis sobre comportamientos \'eticos en la educaci\'on virtual. \emph{Revista Interamericana de Investigaci\'on, Educaci\'on y Pedagog\'{\i}a (RIIEP)}, 14(2), 113--140. \url{https://doi.org/10.15332/25005421.6059}

\bibitem[O'Neil(2017)]{ONeil2017}
O'Neil, C. (2017). \emph{Armas de destrucci\'on matem\'atica: C\'omo el big data aumenta la desigualdad y amenaza la democracia}. Madrid: Capit\'an Swing.

\bibitem[Cortina(1986)]{Cortina1986}
Cortina, A. (1986). \emph{\'Etica m\'{\i}nima: Introducci\'on a la filosof\'{\i}a pr\'actica}. Madrid: Tecnos.

\end{thebibliography}

\end{document}
